\documentclass{article}
\usepackage[UTF8]{ctex}
\usepackage{amsmath,amssymb,amsfonts}
\usepackage{enumitem}
\usepackage{geometry}
\geometry{a4paper, margin=1in}
% === 错题本排版优化（不参与题目标识） ===
\setlength{\parindent}{0pt}
\setlength{\parskip}{0.35em plus 0.10em minus 0.08em}
\setlength{\emergencystretch}{3em}
\setlist{itemsep=0.25em, topsep=0.40em, parsep=0pt}
\allowdisplaybreaks[2]
% === 错题本排版优化结束 ===
\begin{document}
\section*{电子科技大学2024年数学分析真题}
\section*{167．电子科技大学2024年数学分析试卷}
\begin{itemize}

\subsection{题目1}
  \item[1．] 计算 \(\lim _{x \rightarrow 0} \frac{1-(\cos x)^{\sin x}}{\sqrt{1+x^{3}}-1}=\) \(\_\_\_\_\)．

\subsection{题目2}
  \item[2．] 设 \(u(x, y, z)=x y z\) ，则 \(\mathrm{d}^{3} u=\) \(\_\_\_\_\)．
\end{itemize}

\subsection{题目3}
3．曲线积分 \(\oint_{L}\left(2 x^{2}+y^{2}+z\right) \mathrm{d} s=\) \(\_\_\_\_\)，其中 \(L\) 为 \(x^{2}+y^{2}+z^{2}=1\) 与 \(x+y+z=0\) 的交线．\\
\(\_\_\_\_\)

\subsection{题目4}
4．函数 \(z=x y\) 在点 \(\_\_\_\_\)处的法线与平面 \(x+3 y+z+9=0\) 垂直．

\subsection{题目5}
5．级数 \(\sum_{n=1} \overline{n \cdot 2^{n}}=\) \(\_\_\_\_\)．

\subsection{题目6}
6．将 \(f(x)=\pi-x, x \in(0, \pi)\) 展开成正弦级数 \(\_\_\_\_\)，并写出和函数 \(\_\_\_\_\)．计算题．每题 10分，共 60 分．

\subsection{题目7}
7．方程组

\[
\left\{\begin{array}{l}
x y+y+y^{2} z=0 \\
x^{y}+y z-z^{2}+5=0
\end{array}\right.
\]

在点 \(P(1,-2,1)\) 附近能否唯一确定隐函数组？

\subsection{题目8}
8．计算曲线积分

\[
\int_{L} \frac{(x-y) \mathrm{d} x+(x+4 y) \mathrm{d} y}{x^{2}+4 y^{2}} .
\]

其中 \(L: x^{2}+y^{2}=1\) ，取逆时针方向。

\subsection{题目9}
9．求 \(x+2 y=1\) 与 \(x^{2}+2 y^{2}+z^{2}=1\) 相交曲线上到原点距离最小的点．

\subsection{题目10}
10．判断数项级数 \(\sum_{n=1}^{\infty}(-1)^{n-1} \frac{\sin ^{2} n}{n}\) 的敛散性，若收敛需要判断是条件收敛还是绝对收敛。

\subsection{题目11}
11．计算积分 \(I(y)=\int_{0}^{+\infty} e^{-x^{2}} \cos 2 x y \mathrm{~d} x\) ，其中 \(y \in \mathbb{R}\) 。

\subsection{题目12}
12．计算二重积分 \(\iint_{D}\left(\sqrt{x^{2}+y^{2}}+y\right) \mathrm{d} x \mathrm{~d} y\) ，其中 \(D: x^{2}+y^{2}=4\) 与 \((x+1)^{2}+y^{2}=1\) 所围成的区域。

\subsection{题目13}
13．已知数列 \(\left\{a_{n}\right\}\) 满足 \(a_{n}=\sin 1+\frac{\sin 2}{2!}+\frac{\sin 3}{3!}+\cdots+\frac{\sin n}{n!}\) ．

\begin{itemize}
  \item[（1）] 证明：\(\left\{a_{n}\right\}\) 有界，但是 \(\left\{a_{n}\right\}\) 不单调．
  \item[（2）] 证明：数列 \(\left\{a_{n}\right\}\) 收敛．
\end{itemize}

\subsection{题目14}
14．证明：函数项级数 \(\sum_{n=1}^{\infty} \frac{x^{2}}{\left(1+x^{2}\right)^{n}}\) 在 \((-\infty,+\infty)\) 上非一致收敛。

\subsection{题目15}
15．证明：反常积分 \(\int_{0}^{+\infty} \frac{x}{1+x^{5} \sin ^{2} x} \mathrm{~d} x\) 收敛．

\subsection{题目16}
16．已知 \(\lambda>0, \beta \in(0,1),\left\{a_{n}\right\}\) 为正数列，且满足

\[
\liminf _{n \rightarrow \infty} n^{\beta}\left(\frac{a_{n}}{a_{n+1}}-1\right)=\lambda
\]

证明：对任意的 \(k>0\) ，有 \(\lim _{n \rightarrow \infty} n^{k} a_{n}=0\) ．

\subsection{题目17}
17．已知函数 \(f(x) \in C^{2}[0,1]\) ，且满足 \(f(0) f(1) \geq 0\) ．证明：

\[
\int_{0}^{1}\left|f^{\prime}(x)\right| \mathrm{d} x \leq 2 \int_{0}^{1}|f(x)| \mathrm{d} x+\int_{0}^{1}\left|f^{\prime \prime}(x)\right| \mathrm{d} x .
\]
\end{document}
